% Texto do resumo, compartilhado por main.tex (artigo, formato SBC Reviews) e monografia.tex (ABNT).
A engenharia de prompt é uma das formas de usar grandes modelos de linguagem no desenvolvimento de software, e os
estudos sobre o tema variam muito em técnicas, modelos, tarefas e métricas. Este trabalho organiza as evidências sobre
técnicas de prompting aplicadas à geração, tradução, refatoração e depuração de código, por meio de uma revisão
rápida da literatura, de uma taxonomia e de uma tabela de apoio à escolha. A busca no arXiv, no IEEE Xplore e na ACM
Digital Library resultou em 2.304 estudos únicos. Desses, 715 tinham a tarefa no título e foram triados por dois
modelos de linguagem, com conferência do autor. No texto completo, o autor conferiu os 419 estudos lidos pelos modelos e incluiu 332. As fichas dos 332 estudos, com 805 comparações e avaliação de qualidade, foram extraídas por IA, e o autor conferiu
uma amostra de 67 estudos, com correções em 18 de 571 blocos, e conferiu no original as 120 afirmações
empíricas do texto. A maior parte da evidência compara técnicas completas com linhas de base,
sem isolar o mecanismo responsável. Nas comparações extraídas, muitos resultados favoráveis envolvem informação externa
verificável, como testes, execução e contexto do repositório, e dependem de mais chamadas ao modelo ou de informação
da solução de referência. Instruções de raciocínio, autocrítica sem execução, exemplos fixos e mudanças de redação
tiveram resultados divididos, sobretudo com modelos pequenos. Menos da metade dos estudos avalia alguma mitigação de
risco. As principais limitações são a exigência da tarefa no título, a exclusão dos estudos só no arXiv e a extração por IA,
conferida pelo autor só numa amostra.
